\section{Introduction}

Recent advancements in large language models have significantly influenced mathematical reasoning and theorem proving in artificial intelligence. Despite notable progress in natural language domains, language models still encounter substantial challenges in formal theorem proving, \eg~using Lean \citep{moura2021lean} and Isabelle \citep{paulson_isabelle_1994}, which requires rigorous derivations satisfying formal specifications of the verification system.
Even advanced models like GPT-4 \citep{OpenAI} struggle with complex formal proofs, underscoring the intricate nature of both the coding and the mathematics involved.
A formal theorem proving model must not only grasp the syntax and semantics of formal systems like the Lean theorem prover but also align abstract mathematical reasoning with precise formal representation.

Language models in formal theorem proving typically employ two strategies: proof-step generation \citep{polu2020generative, jiang2022thor, lample2022hypertree, yang2024leandojo, wu2024lean} and whole-proof generation \citep{jiang2022draft, zhao2023decomposing, wang2023lego}. Proof-step generation predicts each subsequent tactic and verifies it using the formal verifier to obtain updated information about the current tactic state, often utilizing tree search techniques to construct valid proofs. In contrast, whole-proof generation is computationally efficient, which produces an entire proof code based on the theorem statement, requiring less communication budget to coordinate between the prover model and the formal theorem verifier. While \lastwork~\citep{xin2024deepseek} has achieved state-of-the-art results in Lean 4 with whole-proof generation, this paradigm presents its unique challenges. It requires long-horizon sequence prediction without access to intermediate tactic states, and future tactics depend on these hidden results. In Lean's tactic mode, proofs are constructed through a sequence of tactics that transform the proof state. This sequential nature introduces the risk of compounding errors \citep{ross2011reduction}, where a single misinterpretation can lead to significant deviations from a valid proof path. More specifically, the auto-regressive model may have incorrect believes on intermediate tactic states when generating long proofs.

To seamlessly integrate intermediate tactic states in proof-step generation while maintaining the simplicity and computational efficiency of whole-proof generation, we have developed a unified approach in \thiswork. This method combines the strengths of both proof-step and whole-proof generation techniques through a truncate-and-resume mechanism. The process begins with standard whole-proof generation, where the language model completes the proof code following the theorem statement prefix. The Lean prover then verifies this code. If the proof is correct and complete, the procedure terminates. If an error is detected, the code is truncated at the first error message, and any subsequent code is discarded. The successfully generated proof code is then used as a prompt for the generation of next proof segment. To enhance the accuracy of the model's new completions, we append the latest state from the Lean 4 prover as a comment at the end of the prompt. Notably, our method is not restricted to resuming from the last successfully applied tactic. We integrate the truncate-and-resume mechanism into Monte-Carlo tree search \citep[MCTS;][]{coulom2006efficient} in which the truncation points are scheduled by the tree search policy.
In addition, we propose a novel reward-free exploration algorithm for MCTS to address the reward sparsity issue of proof search.
We assign the tree search agent intrinsic motivation, \aka~curiosity \citep{schmidhuber2010formal}, to extensively explore the tactic state space.
These algorithmic modules extend the functionality of our whole-proof generation model to become a flexible tool for interactive theorem proving, which can effectively utilize the proof assistant feedback and generate diverse solution candidates.

\begin{figure}[htb]
   \centering
   \includegraphics[width=0.75\textwidth]{figures/prover15_main4_training.pdf}
   \caption{\textbf{Overall Framework.} \thiswork~is trained through pre-training, supervised fine-tuning, and reinforcement learning. During supervised fine-tuning, the pre-trained model receives an incomplete theorem proof ending with a tactic state comment keyword. The model is trained to predict the content of this tactic state (auxiliary objective) and complete the subsequent proof steps (main objective). In the reinforcement learning stage, given an incomplete theorem proof and ground-truth tactic state from the Lean prover, we roll out the fine-tuned model to generate multiple proof candidates, which are then verified by the Lean prover. The verification results for these candidates are used as binary (0-1) rewards to further optimize the model and enhance its alignment with the formal specifications of the verification system. For model inference, we offer two alternatives: single-pass sampling and Monte-Carlo tree search.}
  \label{fig:overview_training}
\end{figure}

\subsection{Contributions}

We present a comprehensive framework for developing a language model-based formal mathematics prover, integrating several key components: large-scale mathematical pre-training, formal mathematics corpus construction and augmentation, online reinforcement learning from proof assistant feedback, and a tree search methodology for long-term planning in theorem proving. The pre-trained model, supervised fine-tuned model, and reinforcement learning model, along with the code for the Monte-Carlo tree search algorithm, are publicly available for further research and application.

\begin{itemize}
\item \textbf{Pre-Training}: We enhance our base model's capabilities in formal theorem proving and mathematical reasoning by further pre-training on high-quality mathematics and code data, with a focus on formal languages such as Lean, Isabelle, and Metamath.
\item \textbf{Supervised Fine-Tuning}: We improve the Lean 4 code completion dataset by implementing two data augmentation techniques. First, we use DeepSeek-Coder V2 236B \citep{zhu2024deepseek} to annotate natural language chain-of-thought comments alongside Lean 4 code, aligning formal theorem proving with natural language reasoning. Second, we insert intermediate tactic state information within the Lean 4 proof code, enabling our model to leverage compiler feedback effectively. The resulting dataset is then used to fine-tune the pre-trained model.
\item \textbf{Reinforcement Learning}: We employ the GRPO algorithm \citep{shao2024deepseekmath} to perform reinforcement learning from proof assistant feedback (RLPAF) on the supervised fine-tuned model. Verification results from the Lean prover serve as reward supervision, enhancing the model's alignment with the formal specifications of the verification system.
\item \textbf{Monte-Carlo Tree Search}: We advance the tree search method in formal theorem proving by introducing a novel abstraction and a corresponding search algorithm. Our truncate-and-resume mechanism acts as a state-action abstraction, seamlessly integrating the tree search process into the whole-proof generation framework. We present \treesearch, an innovative Monte-Carlo tree search algorithm that leverages the RMax \citep{brafman2002r} strategy to tackle exploration challenges in sparse-reward proof search problems. By assigning intrinsic rewards, this algorithm encourages the prover agent to generate diverse planning paths, thereby fostering extensive exploration of the proof space.
\end{itemize}

\subsection{Summary of Evaluations and Metrics}

\begin{itemize}
\item \textbf{miniF2F}: In the single-pass whole-proof generation setting, \thiswork~achieved a pass rate of \textbf{60.2\%} on the test set of miniF2F, marking a significant improvement of absolute 10.2 percentage points over \lastwork's 50.0\%. Incorporating tree search techniques further elevated the pass rate to a new state-of-the-art \textbf{63.5\%}.
\item \textbf{ProofNet}: \thiswork~also demonstrated strong performance in the single-pass whole-proof generation setting for ProofNet, with pass rates of \textbf{21.6\%} on the validation set and \textbf{23.7\%} on the test set. The integration of tree search techniques further enhanced these results, achieving new state-of-the-art pass rates of \textbf{25.4\%} on the validation set and \textbf{25.3\%} on the test set.
\end{itemize}

